\subsection{系统级优化方案的设计与实现}

针对上述内存访问瓶颈，本课题提出了四项相互配合的系统级优化技术，并将其整合进\,FastTT\,实现中。
图\ref{fig:opt-overview}给出了\,FastTT\,的整体优化框架：编码阶段先按缓存友好窗口划分（Cache-Friendly Window Partitioning）将数据块切分为小窗口，并以相邻合并访问（Neighbor-Merging Access）模式计算移位异或结果写入校验块；
解码阶段先通过统一解码抽象（Erasure-Parity Mapping）统一数据块与校验块的丢失情况，再依据丢失情况进行遍历模式选择（Traversal Pattern Selection），最后同样以窗口划分方式完成代入与Two-tone消除，一并完成丢失数据块的解码与丢失校验块的修复。

\begin{figure}[htbp]
    \centering
    \includegraphics[width=0.6\textwidth]{Two-tone Optimization Overview.pdf}
    \caption{FastTT的整体优化框架}
    \label{fig:opt-overview}
\end{figure}

\FloatBarrier
\subsubsection{统一解码抽象：Erasure-Parity Mapping}

在通用的存储场景中，解码必须同时保证数据块和校验块的及时恢复。
直观实现需要在解码完成后单独执行一个“修复（Repair）”阶段来重新编码丢失的校验块，引入了冗余操作和额外的内存访问。
为此，本课题提出Erasure-Parity Mapping统一抽象：将所有丢失的数据块和丢失的校验块统一表示为同一种映射结构，每个丢失块按编号顺序匹配一个未丢失的校验块；若丢失块本身就是校验块，则将其映射至自身。
在代入步骤中无需额外判断，直接将未丢失数据块代入所有映射结构；
在消除步骤中，每解出一个码字便立即代入到所有相关映射结构（包括对应丢失校验块的结构）中。
由于数据块的解码与校验块的重编码被统一为对同类映射结构的操作，二者得以相互穿插、同时完成，从而省去了独立的修复阶段，进一步减少了遍历与内存访问次数。
图\ref{fig:naive-flow}与图\ref{fig:mapping-flow}对比了一般解码流程与本方案的统一解码流程。
两图右侧标注了各阶段的复杂度，其中\,MEM\,表示内存访问操作的次数，XOR\,表示异或操作的次数。
对比可见，两种流程的异或操作次数总体相当——这是因为统一解码并未改变底层的Two-tone消除算法逻辑，求解每个丢失块所需的异或运算量是固定的；
真正的差异体现在内存访问上：一般流程需要额外的校验块复制（图\ref{fig:naive-flow}中的“Copy parity chunks”）以及独立的校验块重编码（“Re-encode erased parity chunks”）阶段，二者都引入了大量额外的\,MEM\,操作；
而统一解码流程通过Erasure-Parity Mapping将校验块修复融入消除步骤，省去了上述独立阶段，从而在不增加异或运算的前提下显著减少了内存访问的总次数。

\begin{figure}[htbp]
    \centering
    \begin{minipage}[t]{0.46\textwidth}
        \vspace{0pt}
        \centering
        \includegraphics[width=\linewidth,height=0.34\textheight,keepaspectratio]{Two-tone Naive Decoding Flow.pdf}
        \caption{Two-tone一般解码实现流程}
        \label{fig:naive-flow}
    \end{minipage}%
    \hspace{0.04\textwidth}
    \begin{minipage}[t]{0.46\textwidth}
        \vspace{0pt}
        \centering
        \includegraphics[width=\linewidth,height=0.34\textheight,keepaspectratio]{Erasure-Parity Mapping Decoding Flow.pdf}
        \caption{Erasure-Parity Mapping统一解码流程}
        \label{fig:mapping-flow}
    \end{minipage}
\end{figure}

图\ref{fig:epm-demo}给出了一个具体的解码示例。假设$\boldsymbol d_1$、$\boldsymbol d_2$和$\boldsymbol c_2$丢失，
映射步骤建立三组关联：$\boldsymbol d_1\leftrightarrow\boldsymbol c_3$、$\boldsymbol c_2\leftrightarrow\boldsymbol c_2$（自映射）、$\boldsymbol d_2\leftrightarrow\boldsymbol c_1$；
随后将未丢失的数据块$\boldsymbol d_3$、$\boldsymbol d_4$代入这些映射；在消除过程中，解出的码字按序传播至其他映射结构，丢失校验块$\boldsymbol c_2$的修复也被自然融入消除步骤，从而显著减少冗余内存访问。

\begin{figure}[htbp]
    \centering
    \includegraphics[width=0.62\textwidth]{Erasure-Parity Mapping Decoding Demonstration.pdf}
    \caption{Erasure-Parity Mapping解码举例（假设$\boldsymbol d_1$、$\boldsymbol d_2$、$\boldsymbol c_2$丢失，虚线箭头与删除线标注表示被消除的冗余操作）}
    \label{fig:epm-demo}
\end{figure}

\FloatBarrier
算法\ref{algo:epm-decode}给出了基于Erasure-Parity Mapping抽象的解码过程，包含映射、代入与Two-tone消除三个步骤。
映射步骤将所有丢失的数据块与校验块统一表示为结构\texttt{mpe}，记录每个丢失块及其匹配的存活校验块；若丢失块本身为校验块，则映射至自身。
代入与消除步骤随后对所有\texttt{mpe}条目按统一规则处理，从而省去独立的修复阶段。

\begin{algorithm}[htbp]
\small
\caption{\textsc{Erasure--Parity Mapping Decoding}}
\label{algo:epm-decode}
\KwIn{Ascendingly sorted indices of survived parity chunks $I_{svd}$, indices of erased parity chunks $I_{erd}$, indices of survived data chunks $J_{svd}$, indices of erased data chunks $J_{erd}$, survived data chunks $\boldsymbol{d}_j,j\in J_{svd}$, survived parity chunks $\boldsymbol{c}_i, i\in I_{svd}$, Two-tone offset matrix $(s_{i,j})$}
\KwOut{Decoded data chunks $\boldsymbol{d}_j,j\in J_{erd}$, repaired parity chunks $\boldsymbol{c}_i,i\in I_{erd}$}
\tcc{Line 1-10: $MapErasure(J_{erd}|I_{erd}\leftrightarrow I_{svd})$}
$mpe \gets allocate(m,nullptr)$; $I_{mpd}\gets vector()$\;
\ForEach{$j \in J_{erd}$}{
    $i \gets I_{svd}.pop\_front()$; $I_{mpd}.push\_back(i)$\;
    $mpe[i] \gets \boldsymbol d_j$; \tcp{$d_j[x]$ is undefined}
    $memcpy(\boldsymbol c_i \rightarrow mpe[i], m \times l)$\;
}
\ForEach{$i \in I_{erd}$}{
    $mpe[i] \gets \boldsymbol c_i$; \tcp{$c_i[x]$ is undefined}
    $memset(0 \rightarrow mpe[i], m \times l)$\;
}
$Substitution(\boldsymbol d[J_{svd}] \rightarrow mpe[I_{mpd}+I_{erd}])$\;
$Elimination(mpe[I_{mpd}] \leftrightarrow mpe[I_{mpd}+I_{erd}])$\;
\end{algorithm}

\FloatBarrier
\subsubsection{缓存友好的窗口划分：Cache-Friendly Window Partitioning}

由于移位和异或操作本身极为高效，整体性能受限于数据块在内存中的读写效率，频繁的跨块内存访问极易导致缓存失效。
为此，本课题将每个较长的数据块划分为连续、定长的窗口（window），并在处理下一个窗口前完成当前窗口内的全部移位异或操作，如图\ref{fig:cfw}所示，从而显著提升缓存命中率。
同时，考虑到Ceph等系统分配的块内存通常未预先初始化，本方案采用“首次访问时初始化”（Init On First Access）策略，仅在窗口被实际访问时才进行初始化，避免了对全部块一次性初始化造成的开销与浪费。

\begin{figure}[htbp]
    \centering
    \includegraphics[width=0.6\textwidth]{Cache-Friendly Window Demonstration.pdf}
    \caption{缓存友好窗口划分效果演示。同色片段表示同一窗口内的局部访问，校验块中的阴影区域表示与相邻窗口的重叠}
    \label{fig:cfw}
\end{figure}

在解码过程中，由于代入与消除两个步骤在不同块之间存在错位，若仅对当前窗口代入，可能导致消除位置的数据未被完全覆盖。
为此本方案采用off-by-one代入策略：代入窗口始终比消除窗口领先一个窗口，从而保证代入覆盖消除将要读取的全部位置。
该策略要求窗口大小至少能跨越$k-1$个码字位置，以覆盖Two-tone偏移矩阵引入的全部相对偏移；
据此选取满足$W\ge k-1$的最小2的幂作为窗口大小$W$，兼顾对齐简洁性与覆盖要求，实验证明该取值能取得最佳的吞吐折中。

\FloatBarrier
\subsubsection{内存友好的相邻合并访问：Neighbor-Merging Access}

编码时存在横向与纵向两种遍历模式（如图\ref{fig:access-pattern}(a)、(b)所示）：横向遍历逐个读取数据块码字并立即更新所有对应校验块，数据访问局部性好，但校验块写入分散且重复，内存访问复杂度约为$\Theta((k+km)l)$；
纵向遍历以校验块码字为中心，一次性写入每个校验块码字，减少了校验块写入，却需重复加载数据块，复杂度约为$\Theta((m+km)l)$。两者均存在内存访问瓶颈。

为此本课题提出相邻合并访问策略：在横向遍历的基础上，将相邻的两个数据块码字先加载到SIMD寄存器中，通过一次寄存器内异或合并，再一次性写入对应校验块，如图\ref{fig:access-pattern}(c)所示。
相邻合并访问的内存访问复杂度降低为约$\Theta((k+\tfrac{km}{2})l)$，明显优于横向遍历，并在$m(k+2)>2k$（实际$(k,m)$配置普遍满足）时也优于纵向遍历。
该方法既避免了对同一数据块内存的重复访问，又将对校验块内存的访问理论上减少一半，同时通过指令级两块异或提升了SIMD通道吞吐利用率。
考虑到SIMD寄存器数量有限且更多校验块会增加偏移复杂度，本方案有意将合并限制在相邻两块之间。

\begin{figure}[htbp]
    \centering
    \begin{minipage}[b]{0.30\textwidth}
        \centering
        \includegraphics[height=4.3cm,keepaspectratio]{Two-tone Horizontal Access Pattern.pdf}\\[2pt]
        \textbf{(a) 横向遍历}
    \end{minipage}%
    \hfill
    \begin{minipage}[b]{0.66\textwidth}
        \centering
        \includegraphics[height=4.3cm,keepaspectratio]{Two-tone Vertical Access Pattern.pdf}\\[2pt]
        \textbf{(b) 纵向遍历}
    \end{minipage}

    \vspace{8pt}

    \begin{minipage}[b]{\textwidth}
        \centering
        \includegraphics[height=4.8cm,keepaspectratio]{Neighbor-Merging Access Pattern.pdf}\\[2pt]
        \textbf{(c) 相邻合并访问}
    \end{minipage}
    \caption{Two-tone编码的访问模式对比：(a)横向遍历，数据局部性好但校验写入分散；(b)纵向遍历，校验写入集中但重复读取数据；(c)相邻合并访问，相邻数据块窗口先在SIMD寄存器内异或合并，再以一次连续存储写回各校验块窗口}
    \label{fig:access-pattern}
\end{figure}

\FloatBarrier
\subsubsection{遍历模式选择：Traversal Pattern Selection}

在大规模存储集群中，故障模式高度偏向单块丢失，同时丢失两块及以上的概率远低于丢失一块。
为此本课题针对“仅丢失一个数据块或一个校验块”的常见场景设计了特化优化。
此时解码无需执行完整的Two-tone消除流程，只需对所有相关的未丢失块做一次全局异或，一次性恢复丢失块。
由于该操作对每个相关码字恰好访问一次，特别适合采用纵向遍历模式（此时无重复访问，效率高于横向模式），并可进一步采用AVX2的stream类指令让内存访问绕过缓存，省去缓存相关处理流程，如图\ref{fig:special}所示。
该特化分支与统一解码框架完美兼容：解码器先执行Erasure-Parity Mapping构造映射，再根据丢失块数量选择遍历模式——丢失一块时走单遍纵向异或快路径，否则回退到标准的“代入+消除”流程。

\begin{figure}[htbp]
    \centering
    \includegraphics[width=0.6\textwidth]{Vertical Traversal Pattern Demonstration.pdf}
    \caption{单块丢失场景下的纵向遍历模式演示：丢失块（如$\boldsymbol d_2$或$\boldsymbol c_1$）通过对未丢失块的一遍全局异或一次性恢复，无需Two-tone消除}
    \label{fig:special}
\end{figure}

算法\ref{algo:tps-decode}给出了带遍历模式选择的统一解码过程。
解码器先执行Erasure-Parity Mapping构造映射结构，再依据丢失块总数选择遍历模式：当恰好丢失一个块（数据块或校验块）时，跳过完整的Two-tone消除，仅以一遍纵向遍历对所有必要的存活块做全局异或、一次性重建丢失块；否则回退到标准的“代入+Two-tone消除”流程。

\begin{algorithm}[htbp]
\small
\caption{\textsc{Traversal Pattern Selection Decoding}}
\label{algo:tps-decode}
\KwIn{Ascendingly sorted indices of survived parity chunks $I_{svd}$, indices of erased parity chunks $I_{erd}$, indices of survived data chunks $J_{svd}$, indices of erased data chunks $J_{erd}$, survived data chunks $\boldsymbol{d}_j,j\in J_{svd}$, survived parity chunks $\boldsymbol{c}_i, i\in I_{svd}$, Two-tone offset matrix $(s_{i,j})$}
\KwOut{Decoded data chunks $\boldsymbol{d}_j,j\in J_{erd}$, repaired parity chunks $\boldsymbol{c}_i,i\in I_{erd}$}
$MapErasure(J_{erd}\,|\,I_{erd}\leftrightarrow I_{svd})$\;
\tcc{Traversal Pattern Selection}
\If{$|J_{erd}|+|I_{erd}|==1$} {
    \tcc{Vertical pattern reconstruction (single pass)}
    $i\gets I_{mpd}[1]$\;
    \For{$x \gets 1$ \KwTo $l$} {
        $sum\gets 0$\;
        \ForEach{$j \in J_{svd}$} {
            $sum\gets sum \oplus d_j[x-s_{ij}]$\;
        }
        $mpe[i][x] \gets sum$\;
    }
}
\Else {
    \tcc{Standard unified decoding}
    $Substitution(\boldsymbol d[J_{svd}] \rightarrow mpe[I_{mpd}+I_{erd}])$\;
    $Elimination(mpe[I_{mpd}]~from~mpe[I_{mpd}+I_{erd}])$\;
}
\end{algorithm}

\FloatBarrier
